\documentclass[11pt,a4paper,sans]{moderncv}

% moderncv themes
\moderncvstyle{classic}
\moderncvcolor{blue}

\usepackage[scale=0.75]{geometry}
\setlength{\footskip}{149.60005pt}

\ifxetexorluatex
  \usepackage{fontspec}
  \usepackage{unicode-math}
  \defaultfontfeatures{Ligatures=TeX}
  \setmainfont{Latin Modern Roman}
  \setsansfont{Latin Modern Sans}
  \setmonofont{Latin Modern Mono}
  \setmathfont{Latin Modern Math}
\else
  \usepackage[utf8]{inputenc}
  \usepackage[T1]{fontenc}
  \usepackage{lmodern}
\fi

\usepackage[english]{babel}
\usepackage{multibib}

% personal data
\name{Matteo}{Iervasi}
\email{matteoiervasi@gmail.com}
\homepage{matteoiervasi.it}

% Social icons
\social[linkedin]{matteo-iervasi}
\social[orcid]{0009-0004-9029-9479}
\social[googlescholar][scholar.google.com/citations?user=2DQ1LC0AAAAJ\&hl=it]{Google Scholar}

\photo[64pt][0.4pt]{picture}
\renewcommand*{\bibliographyitemlabel}{[\arabic{enumiv}]}
\begin{document}
\makecvtitle

I'm a firmware engineer with a strong background in low-level programming and microcontroller development, experienced in creating efficient and 
reliable firmware solutions for embedded systems and constantly seeking opportunities to apply my technical expertise to contribute to innovative 
and impactful projects. I'm an expert in C, C++ and Python, especially in optimizing performance and ensuring hardware compatibility and I'm also
highly skilled in utilizing the Yocto Project to customize embedded Linux distributions.

\section{Education}
\cventry{2024--Present}{Ph.D. fellowship in AI for Active and Assisted Living}{University of Stavanger}{Stavanger}{}{
	Research on the development of AI-based solutions for the improvement of the quality of life of Parkinson's patients.
}
\cventry{2018--2020}{Master's degree in Computer Science and Engineering}{University of Verona}{Verona}{}{
	Thesis title: "Integrating synthetic and real components of a cyber-physical production system".\\
	Relevant courses:
	\begin{itemize}
		\item Embedded systems design
		\item Networked embedded systems
		\item Physics of integrated devices
		\item System theory
	\end{itemize}
}
\cventry{2015--2018}{Bachelor's degree in Computer Science}{University of Verona}{Verona}{}{
	Relevant courses:
	\begin{itemize}
		\item Operating systems
		\item Software engineering
		\item Signal and image processing
		\item Language and compilers
	\end{itemize}
}
\cventry{2010--2015}{High school diploma}{Liceo Scientifico Angelo Messedaglia}{Verona}{}{
	Attended the ``\textit{Applied Science}'' curriculum, which focuses on Computer Science, Physics, Chemistry and Biology.
	\begin{itemize}
		\item Developed a software for controlling the chemistry lab spectrophotometer
	\end{itemize}
}

\section{Master thesis}
\cvitem{Title}{\emph{Integrating synthetic and real components of a cyber-physical production system}}
\cvitem{Supervisor}{Prof. Franco Fummi}
\cvitem{Advisor}{Dr. Stefano Centomo}
\cvitem{Abstract}{
	One of the key aspects of \emph{Industry 4.0} is the concept of \texttt{Digital Twins}, as they are an enabling
	technology for things like predictive maintenance, real-time production optimization, on-demand product
	customization and so on\ldots
	A limiting factor in the creation of \emph{Digital Twins} is the abundance of incompatible modeling languages. Among
	the research and the projects that tries to overcome this issue, \texttt{AutomationML} is increasingly cited and used
	as a vendor-neutral language for model exchange.
	This work proposes a simple direct approach for the integration of models in CPPS systems, using \texttt{AutomationML}
	as the base technology.
}

\section{Internships}
\cventry{2021}{Research internship}{EDALab S.r.l.}{San Giovanni Lupatoto}{}{Integration of Bluetooth Low Energy devices in BOX-IO system.\\
Focus of the work was the development of a firmware for some microcontrollers (most notably Texas Instruments and ST) able to communicate over 
BLE, ZigBee and generic serial connection with the BOX-IO system.}

\section{Employment}
\cventry{Oct 2020-- Jul 2024}{Embedded software engineer}{EDALab S.r.l.}{San Giovanni Lupatoto}{}{General firmware and low-level embedded software development for third-parties.
\begin{itemize}
	\item Collaborated directly with customers to design and develop firmware solutions for various embedded systems.
	\item Implemented low-level software and optimized performance while ensuring hardware compatibility.
	\item Contributed to the development of real-time systems for industrial applications.
\end{itemize}
At EDALab, I had the opportunity to work for over 30 projects across diverse applications. I programmed firmwares that now power HVAC 
plants, industrial vapor heaters, thermostats, compressors for oil extraction, and even coffee machines, along with other mixed projects
that include sensor based monitoring activities.
I particularly enjoyed contributing to projects improving people's lives, such as developing firmware for wearable devices powered by Espressif ESP32, 
leveraging my expertise in low-power microcontrollers.
}
\cventry{Jan 2024}{Teaching professor}{University of Verona}{Verona}{}{Conducted a course entitled ``Firmware Development with Bluetooth Low Energy 
	Protocol and FreeRTOS."}
\cventry{Aug 2020}{Internship}{EDALab S.r.l.}{San Giovanni Lupatoto}{}{Integrated reliable update mechanism base on SWUpdate for BoxIO.}
\cventry{Dec 2019}{Internship}{University of Verona}{Verona}{}{Developed an operating system image for classrooms' displays based on Raspberry Pi.}
\cventry{Jan 2017-- Mar 2018}{Internship}{Sordato S.r.l.}{Monteforte d'Alpone}{}{Developed software for controlling an array of automated wine fermentation machines.}
\cventry{2017--2018}{Teacher assistant}{University of Verona}{Verona}{}{I worked as a teacher assistant in the following courses:
\begin{itemize}
	\item Operative systems
	\item Programming I
	\item Programming II
\end{itemize}
}

\subsection{Volunteering}
\cventry{2017-- Present}{Technician}{AVIS}{Vigasio}{}{I volunteered as a technician and general assistance at the local AVIS association,
	an Italian organization that promotes blood donation.}

\section{Languages}
\cvitemwithcomment{Italian}{Native language}{}
\cvitemwithcomment{English}{Professional level. Proficient in reading and writing, fluent in speaking.}{IELTS certification.}

\section{Skills}
\cvitem{Programming Languages}{Excellent knowledge of C, C++ and Python.}
\cvitem{Firmware Development}{Expertise in low-level programming for various microcontrollers with various environments and compilers, e.g. IAR®, GCC/Clang and Keil®.}
\cvitem{MCUs}{Expertise in developing for ST® STM32, Renesas® RL78, Renesas® RX130, Renesas® RA, NXP® Kinetis, Cypress® FM4 and other ARM MCUs, Microchip® PIC18, PIC24 and PIC32, Intel® 8051 and Espressif® ESP32/ESP8266.}
\cvitem{Embedded Systems}{Experienced using the Yocto Project for customizing embedded Linux distributions, Qt/QML for developing complex HMIs.}
\cvitem{Operating Systems}{Proficient in developing on GNU/Linux and Microsoft® Windows®.}
\cvitem{Scripting}{Expert in OS automation using Bash and general automation using Python.}
\cvitem{Versioning}{Excellent knowledge of Git.}


\section{Interests}
\cvitem{Electronics}{I like studying electronics and then putting it into practice, from circuit design to PCB printing.}
\cvitem{3D printing}{I'm into 3D printing and I enjoy designing things in CAD that then I use in my projects.}
\cvitem{DIY}{I love to repair home appliances, furnitures and fixing broken computers.}


\section{Publications}
\renewcommand*{\bibliographyhead}[1]{}
\bibliographystyle{plain}
\nocite{*}
\bibliography{publications}

\clearpage

\end{document}

